\documentclass[12pt,english]{article}
\usepackage{mathptmx}
\usepackage[top=1in, bottom=1in, left=1in, right=1in]{geometry}
\usepackage{amsmath,amstext,amssymb}
\usepackage{setspace}
\usepackage{booktabs}
\usepackage{graphicx}
\usepackage[english]{babel}
\usepackage[unicode=true,pdfusetitle,
 bookmarks=true,colorlinks,citecolor=black,filecolor=black,
 linkcolor=black,urlcolor=black]{hyperref}

% modelsummary / tinytable / datasummary_skim preamble
\usepackage{tabularray}
\usepackage{float}
\usepackage{codehigh}
\usepackage[normalem]{ulem}
\UseTblrLibrary{booktabs}
\UseTblrLibrary{siunitx}
\newcommand{\tinytableTabularrayUnderline}[1]{\underline{#1}}
\newcommand{\tinytableTabularrayStrikeout}[1]{\sout{#1}}
\NewTableCommand{\tinytableDefineColor}[3]{\definecolor{#1}{#2}{#3}}

\title{Problem Set 11\\ Econ 5253 -- Spring 2026}
\author{Yeyang Zhou}
\date{April 27, 2026}

\begin{document}

\maketitle

\section*{Question 6: Summary Statistics}

\input{datasummary.tex}

The summary statistics make sense for a sample of working women in 1988.
Mean years of schooling \texttt{hgc} is around the high-school-plus-some-college
range, and \texttt{exper} averages a moderate number of years on the
current job. Roughly 20\% of the sample is in a union job, around 60\%
are married, and roughly half have at least one child living at home.

The missing rate of \texttt{logwage} is \textbf{30.7\%}. Given the
sample is restricted to women already known to be working, the
missingness is most likely \textbf{Missing Not At Random (MNAR)} or at
best Missing At Random (MAR): wages are most likely missing for women
whose wage outcomes are systematically related to unobserved factors
(e.g., women who work informally or report wages selectively). The
selection problem is exactly what the Heckman model in Q7 is designed
to address.

\section*{Question 7: Imputation Methods and Returns to Schooling}

I estimate
\begin{equation}
\label{eq:wage}
\log wage_i = \beta_0 + \beta_1 hgc_i + \beta_2 union_i
            + \beta_3 college_i + \beta_4 exper_i
            + \beta_5 exper_i^2 + \eta_i
\end{equation}
under three approaches to handling missing \texttt{logwage}:
(i) listwise deletion (complete cases), (ii) mean imputation, and
(iii) Heckman two-step selection.

\input{regtable.tex}

The coefficient of interest $\beta_1$ (returns to schooling, true value
0.091) takes the following values:

\begin{center}
\begin{tabular}{lc}
\toprule
Method & $\hat{\beta}_1$ \\
\midrule
Complete cases & 0.0590 \\
Mean imputation & 0.0363 \\
Heckman 2-step & 0.0915 \\
\midrule
True value & 0.0910 \\
\bottomrule
\end{tabular}
\end{center}

\textbf{Discussion.} Complete-cases regression underestimates the true
returns to schooling at 0.059, because the missing observations are
non-random and listwise deletion drops them in a way that biases the
remaining sample. Mean imputation performs the worst at 0.036; filling
in missing values with the unconditional mean compresses the variance
of $\log wage$ and attenuates the slope on \texttt{hgc} toward zero.
The Heckman two-step estimator recovers $\hat{\beta}_1 = 0.0915$,
essentially the true value. By explicitly modeling the selection
equation (the probability of observing a wage as a function of
\texttt{married} and \texttt{kids}, in addition to the wage covariates),
Heckman corrects for the non-random nature of the missingness and
delivers a consistent estimate.

The pattern confirms the standard ordering of these methods under
non-random missingness: mean imputation is the worst, complete-cases
is biased but not catastrophic, and a properly-specified selection
model is best when an exclusion restriction is available.

\section*{Question 8: Probit Model of Union Job Preference}

The probit model
$$ u_i = \alpha_0 + \alpha_1 hgc_i + \alpha_2 college_i + \alpha_3 exper_i
       + \alpha_4 married_i + \alpha_5 kids_i + \varepsilon_i $$
yields the following estimates:

\begin{center}
\begin{tabular}{lcccc}
\toprule
Variable & Estimate & Std.\ Error & $z$ & $p$ \\
\midrule
(Intercept) & $-$6.743 & 0.804 & $-$8.39 & $<0.001$ \\
\texttt{hgc}     & $-$1.009 & 0.098 & $-$10.34 & $<0.001$ \\
\texttt{college} &    0.397 & 0.427 &    0.93 & 0.352 \\
\texttt{exper}   &    1.849 & 0.156 &   11.86 & $<0.001$ \\
\texttt{married} &    0.588 & 0.206 &    2.86 & 0.004 \\
\texttt{kids}    &    0.799 & 0.202 &    3.96 & $<0.001$ \\
\bottomrule
\end{tabular}
\end{center}

\section*{Question 9: Counterfactual Policy}

I simulate a policy in which there is no preference or penalty to wives
or mothers for working in union jobs by setting the coefficients on
\texttt{married} and \texttt{kids} to zero and recomputing predicted
probabilities. The results are:

\begin{center}
\begin{tabular}{lc}
\toprule
 & Avg.\ $P(\text{union})$ \\
\midrule
Baseline                 & 0.2373 \\
Counterfactual ($\alpha_4 = \alpha_5 = 0$) & 0.2174 \\
\midrule
Difference                                  & $-$0.0199 \\
\bottomrule
\end{tabular}
\end{center}

Removing the marriage and kids effects lowers the average predicted
probability of holding a union job by about 2 percentage points. This
indicates that, in the estimated model, being married and having kids
\emph{raises} the probability of union employment, so removing those
effects mechanically lowers it.

\paragraph{Is the model realistic?}
No. Several features of the estimates raise concerns:

\begin{enumerate}
  \item The coefficient on \texttt{hgc} is $-1.009$, implying that one
        additional year of schooling \emph{decreases} the latent utility
        of union employment by an enormous amount. This is implausibly
        large in magnitude and likely the wrong sign---empirically,
        union membership in the late 1980s was concentrated among
        workers with moderate education, not strongly negatively related
        to schooling at this magnitude.
  \item The coefficient on \texttt{exper} is $+1.849$, which is also
        implausibly large in absolute value for a probit model.
  \item R issued the warning ``fitted probabilities numerically 0 or 1
        occurred,'' which is a classic symptom of separation or near
        separation in the data and a signal that the model is
        numerically unstable.
  \item The model omits demand-side variables (industry, occupation,
        geography) that drive union availability. With those omitted,
        the coefficients on individual characteristics absorb correlated
        omitted-variable effects.
\end{enumerate}

The counterfactual is therefore best read as descriptive: it tells us
how the average predicted probability would change if the
\texttt{married} and \texttt{kids} coefficients were zeroed out
\emph{within this particular fitted model}, not as a credible estimate
of how a real-world policy would change union participation. A more
defensible model would include industry/occupation controls and would
not produce coefficients of magnitude greater than one on continuous
covariates.

\end{document}
